La sicurezza informatica rappresenta oggi una delle sfide più critiche nell'ambito delle comunicazioni e delle reti informatiche. 
Data la continua evoluzione delle tecnologie e dato il costante aumento nella complessità dei sistemi informatici; emergono costantemente nuove vulnerabilità e tecniche di 
attacco che permetteranno l'esfiltrazione dei dati. Ciò ha richiesto un'attenzione particolare diventando un ambito primario di ricerca sia da parte di ricercatori universitari 
che di professionisti del settore. 
Tra queste tecniche, i covert channel costituiscono una minaccia particolarmente insidiosa, in quanto permettono la trasmissione di informazioni (fra due entità) in maniera 
non autorizzata sfruttando protocolli, canali di comunicazione e/o risorse legittimi. 
Ciò rendendo complessa la loro individuazione da parte dei tradizionali sistemi di sicurezza.
\vspace{2ex} \newline  
Il presente lavoro di tesi si concentra sull'analisi e sullo sviluppo di un covert channel basato sul protocollo ICMP (Internet Control Message Protocol) per 
l'esfiltrazione di dati. Quest'ultimo risulta fondamentale per il funzionamento, il monitoraggio e la diagnostica delle reti IP e che, proprio per questo, non può essere 
completamente disabilitato o filtrato nelle reti moderne senza comprometterne l'operatività. 
Tutto questo lo rende un candidato ideale per la creazione di canali di comunicazione difficilmente rilevabili.
L'obiettivo principale è dimostrare come il protocollo ICMP possa essere sfruttato per l'implementazione di diverse tipologie di Covert Channel che sono: 
i timing covert channel, gli  storage covert channel ed i behavioural covert channel.  
\vspace{2ex} \newline 
La tesi si articolerà in sei capitoli. 
\begin{itemize} 
    \item Nel Capitolo 2 verranno inquadrati i Covert Channel ed il protocollo ICMP descrivendone i relativi concetti fondamentali; analizzando le diverse tipologie e 
    caratteristiche dei canali nascosti sviluppabili e alle diverse tipologie di messaggi ICMP disponibili.  
    \item Il Capitolo 3 descriverà gli strumenti software utilizzati per lo sviluppo del progetto ed il relativo testing. 
    In esso sarà presente anche un analisi di progetti esistenti che hanno già affrontato problematiche simili, identificandone punti di forza e criticità. 
    Le osservazioni ricavate hanno poi guidato la progettazione dell'architettura descritta nel capitolo successivo. 
    \item Il Capitolo 4 costituisce la parte centrale dell'elaborato: presenta l'implementazione dettagliata dell'architettura del sistema e della comunicazione fra le 
    entità coinvolte (nello scambio delle informazioni). Descrive inoltre l'implementazione e lo sviluppo delle diverse tipologie di covert channel presenti. 
    Di particolare importanza saranno le metodologie con cui si sono nascosti i dati sia attraverso i tempi dei messaggi sia attraverso i campi presenti in un pacchetto. 
    \item I risultati dei test effettuati verranno riportati nel Capitolo 5. Ciò verrà fatto utilizzando RITA, come strumento per il rilevamento delle intrusioni, 
    ed analizzando come diverse variazioni in alcuni parametri del covert channel influenzino la probabilità di essere individuati. 
    In particola i test potranno variare in base alla quantità di dati trasmessi, alla frequenza di invio dei dati, l'utilizzo dei mittenti reali o falsificati e la presenza 
    (o la mancanza) di ritardi tra una trasmissione e la successiva.  
    \item Infine, il Capitolo 6 presenta le conclusioni del lavoro e riassume i risultati ottenuti, evidenziando i fattori che influenzano maggiormente la rilevabilità di 
    un covert channel. 
    Vengono inoltre fornite raccomandazioni per il monitoraggio e limitazione del traffico ICMP nelle reti, al fine di mitigare i rischi associati a questo tipo di minaccia. 
\end{itemize}  
%Questo lavoro contribuisce alla comprensione delle vulnerabilità intrinseche del protocollo ICMP e fornisce una base per lo sviluppo di contromisure più efficaci contro i 
%covert channel basati su questo protocollo. La ricerca sottolinea l'importanza di un approccio proattivo alla sicurezza di rete, che non si limiti al blocco di traffico 
%manifestamente dannoso, ma che includa anche l'analisi comportamentale e statistica del traffico apparentemente legittimo. 

